% Options for packages loaded elsewhere
\PassOptionsToPackage{unicode}{hyperref}
\PassOptionsToPackage{hyphens}{url}
%
\documentclass[
]{article}
\usepackage{amsmath,amssymb}
\usepackage{lmodern}
\usepackage{ifxetex,ifluatex}
\ifnum 0\ifxetex 1\fi\ifluatex 1\fi=0 % if pdftex
  \usepackage[T1]{fontenc}
  \usepackage[utf8]{inputenc}
  \usepackage{textcomp} % provide euro and other symbols
\else % if luatex or xetex
  \usepackage{unicode-math}
  \defaultfontfeatures{Scale=MatchLowercase}
  \defaultfontfeatures[\rmfamily]{Ligatures=TeX,Scale=1}
\fi
% Use upquote if available, for straight quotes in verbatim environments
\IfFileExists{upquote.sty}{\usepackage{upquote}}{}
\IfFileExists{microtype.sty}{% use microtype if available
  \usepackage[]{microtype}
  \UseMicrotypeSet[protrusion]{basicmath} % disable protrusion for tt fonts
}{}
\makeatletter
\@ifundefined{KOMAClassName}{% if non-KOMA class
  \IfFileExists{parskip.sty}{%
    \usepackage{parskip}
  }{% else
    \setlength{\parindent}{0pt}
    \setlength{\parskip}{6pt plus 2pt minus 1pt}}
}{% if KOMA class
  \KOMAoptions{parskip=half}}
\makeatother
\usepackage{xcolor}
\IfFileExists{xurl.sty}{\usepackage{xurl}}{} % add URL line breaks if available
\IfFileExists{bookmark.sty}{\usepackage{bookmark}}{\usepackage{hyperref}}
\hypersetup{
  pdftitle={Doing things right vs doing right things},
  pdfauthor={Kishore Puthezhath},
  hidelinks,
  pdfcreator={LaTeX via pandoc}}
\urlstyle{same} % disable monospaced font for URLs
\usepackage[margin=1in]{geometry}
\usepackage{graphicx}
\makeatletter
\def\maxwidth{\ifdim\Gin@nat@width>\linewidth\linewidth\else\Gin@nat@width\fi}
\def\maxheight{\ifdim\Gin@nat@height>\textheight\textheight\else\Gin@nat@height\fi}
\makeatother
% Scale images if necessary, so that they will not overflow the page
% margins by default, and it is still possible to overwrite the defaults
% using explicit options in \includegraphics[width, height, ...]{}
\setkeys{Gin}{width=\maxwidth,height=\maxheight,keepaspectratio}
% Set default figure placement to htbp
\makeatletter
\def\fps@figure{htbp}
\makeatother
\setlength{\emergencystretch}{3em} % prevent overfull lines
\providecommand{\tightlist}{%
  \setlength{\itemsep}{0pt}\setlength{\parskip}{0pt}}
\setcounter{secnumdepth}{-\maxdimen} % remove section numbering
\ifluatex
  \usepackage{selnolig}  % disable illegal ligatures
\fi

\title{Doing things right vs doing right things}
\author{Kishore Puthezhath}
\date{28-02-2021}

\begin{document}
\maketitle

{
\setcounter{tocdepth}{2}
\tableofcontents
}
\hypertarget{the-story-so-far}{%
\subsubsection{The story so far:}\label{the-story-so-far}}

\includegraphics{/https://nascenia.com/are-you-a-leader-or-a-boss/}

In the last week of December 2020, I woke up suddenly being added as a
member in yet another WhatsApp group:Medical Education Unit. Even the
director, Academic coordinator and the principal were members. Though
having experienced disappointment a lot in the past, I felt some
excitement. When the Principal posted that the purpose is to ``empower''
the faculty, my excitement grew and I started researching and dreaming
about a few things.

\begin{itemize}
\tightlist
\item
  What the the good \texttt{medical\ education\ Journal} available in
  the corner?
\item
  How to improve the academic atmosphere?
\item
  Are we going to get a great state of art medical education facility?
\item
  Are we going to discuss about the simulation lab?
\end{itemize}

\hypertarget{what-are-the-key-challenges}{%
\subsubsection{What are the key
challenges?}\label{what-are-the-key-challenges}}

The pandemic raises serious questions for medical colleges in India.
During this unprecedented time, as trainee doctors, medical students
stand to learn a tremendous amount and can contribute to the care of
patients. How ever, clinical rotations can expose medical students to
COVID-19 positive patients, risking their safety. Moreover, infected
students may themselves become the source of spread of the disease.
Physical distancing measures have precluded the conduct of
classroom-based teaching and discussions, even for preclinical students.

\begin{itemize}
\tightlist
\item
  \texttt{The\ day-by-day\ time\ table\ for\ new\ CBME\ (Competency-based\ Medical\ Education)\ batch\ of\ 2019\ is\ disrupted}.
\item
  Additionally, \textbf{it was about the time to prepare for phase 2
  timetable}.
\item
  Teachers of phase 2 also needed training for the new curriculum, which
  has been paused.
\item
  Teachers and postgraduate students are being \texttt{recruited} for
  various COVID19 activities. Even \texttt{Professors} are doing
  clerical works like \texttt{records\ keeping} and some of them are
  doing \texttt{swabbing}.
\end{itemize}

In effect time available for teaching has reduced. More over, medical
students will suffer a reduced exposure to certain clinical branches or
a proportionally shortened rotation in all clinical branches. In case
the pandemic continues for an unforeseeable time, an extension of the
medical training period may be warranted. In all situations, the medical
students will be at a loss.

\hypertarget{have-we-done-anything-to-help-them}{%
\subsubsection{Have we done anything to help
them?}\label{have-we-done-anything-to-help-them}}

\textbf{At most institutions, students have continued the curriculum but
in a purely ``non‐clinical'' format}.

Faculty have continued to lecture using prerecorded videos or online
video conferencing platforms like \texttt{WebEx,\ Zoom}, and so forth.
Medical universities such as Kerala University of Health Sciences have
come up with guidelines and limited training programs for the faculty on
the usage of online platforms. In a move triggered by the COVID-19
pandemic, the board of governors(BoG) in supersession of the Medical
Council of India (MCI) recently introduced a new pandemic management
module to better prepare the future doctors for challenges posed by
pandemics like COVID-19.

\hypertarget{why-should-medical-education-continue}{%
\subsection{Why should medical education
continue?}\label{why-should-medical-education-continue}}

There are several important reasons for ensuring the continuation of
medical education in this hour.
\texttt{This\ pandemic\ could\ be\ a\ once\ in\ a\ lifetime\ experience\ for\ anybody,\ not\ the\ least\ for\ a\ future\ medical\ practitioner}.
If the current pandemic is going to take a reasonable time to abate,
there may arise a paucity of healthcare workers. In such situations,
they may have to engage in certain aspects of patient care. If clinical
rotations are deferred for the current student batches and clubbed with
others, the density of learners would impair the clinical learning
experience. A student who misses any part of the medical education is
likely to find it difficult to join the dots later. Therefore, we must
continue to facilitate the ongoing knowledge and skill development of
the next generation of health professionals.

\hypertarget{is-remote-learning-the-solution}{%
\subsection{Is remote learning the
solution?}\label{is-remote-learning-the-solution}}

Despite the undoubted advantage of e-learning and simulation-based
education, in the current situation, there are distinct disadvantages.
\textbf{These tools can be at the best supplemental to clinical teaching
but not a replacement}. Moreover, setting up of a virtual learning or
simulation environment is costly and time taking. Another major
challenge of remote learning is the digital divide and the disparity in
access -- from electricity and internet connections to devices like
computer or smartphones.

\hypertarget{what-is-the-way-forward}{%
\subsection{What is the way forward?}\label{what-is-the-way-forward}}

We need to prepare both the medical students and the faculty for
continuing medical education, not only for now but also for possible
future contagions. It will be reasonable to invest in these pedagogical
innovations over the long-term and develop a central repository of the
requisite sources beforehand. It may also be the right time for
regulators, universities, educational experts and professional
associations to come together as a team and lead the way forward in
these difficult times.
\texttt{We\ have\ to\ seriously\ rethink\ the\ strategy\ of\ keeping\ the\ medical\ students\ away\ from\ patients.}

May this period of `no teaching' be the period of `greatest learning.'

\texttt{edit}: Students started coming to the hospitals 1 week back.

\end{document}
